

\chapter{Background}
\label{chap:background}

\section{MITRE ATT\&CK}

MITRE ATT\&CK es el marco de referencia fundamental para describir y clasificar el comportamiento de actores maliciosos en ejercicios de \textit{Adversary Emulation}. Su estructura proporciona un vocabulario común basado en tres niveles de abstracción:

\begin{itemize}
    \item \textbf{Tácticas (el ``para qué'')}: Representan los objetivos tácticos del adversario en cada fase del ataque. Describen la meta que se pretende alcanzar, como obtener acceso inicial, elevar privilegios o exfiltrar información.

    \item \textbf{Técnicas (el ``qué'')}: Definen el método concreto que utiliza el adversario para cumplir una táctica. Representan la acción observable destinada a lograr el objetivo, como el uso de \textit{phishing}, ejecución de comandos o inyección de procesos.

    \item \textbf{Procedimientos (el ``cómo'')}: Son la implementación específica de una técnica. Detallan los pasos exactos, herramientas, parámetros y artefactos utilizados por un atacante. Un procedimiento debe ser lo suficientemente preciso como para que el \textit{Blue Team} pueda generar reglas de detección y el \textit{Red Team} pueda reproducirlo fielmente.
\end{itemize}

Este modelo estructurado permite mapear comportamientos reales, estandarizar evaluaciones y alinear las actividades defensivas y ofensivas dentro de un lenguaje común.

\section{Cyber Kill Chain (CKC)}

La \textit{Cyber Kill Chain}, desarrollada por Lockheed Martin, describe las fases operativas de un ciberataque con el objetivo de estructurar y contextualizar las acciones del adversario. No se trata de un marco rígido, sino de un modelo conceptual que divide el ataque en etapas bien definidas para facilitar su análisis, detección e interrupción. El modelo identifica las tareas que un adversario debe completar para alcanzar su objetivo final, permitiendo anticipar comportamientos y reforzar las defensas en cada fase.

Los siete pasos de la \textit{Cyber Kill Chain} mejoran la visibilidad del analista sobre el ataque y enriquecen su comprensión de las tácticas, técnicas y procedimientos (TTPs) empleados por el adversario.

\begin{enumerate}
    \item \textbf{Reconocimiento}: Investigación del objetivo y recolección de información.
    \item \textbf{Armamento}: Preparación o creación del payload y capacidades ofensivas asociadas.
    \item \textbf{Entrega}: Transmisión del payload mediante el vector seleccionado.
    \item \textbf{Explotación}: Activación del payload aprovechando una vulnerabilidad o interacción del usuario.
    \item \textbf{Instalación}: Establecimiento del implante y preparación del entorno persistente.
    \item \textbf{Mando y Control (C2)}: Creación del canal de comunicación entre el host comprometido y el adversario.
    \item \textbf{Acciones sobre Objetivos}: Ejecución de los objetivos finales, como movimiento lateral, exfiltración o impacto.
\end{enumerate}

\section{Threat Intelligence}

La Inteligencia de Amenazas es esencial para la simulación de adversarios (\textit{Adversary Simulation}), ya que transforma datos sobre actores maliciosos, campañas y vulnerabilidades en información procesable. Se nutre de análisis técnicos, alertas de organismos de ciberseguridad, investigaciones de la industria y colecciones especializadas de TTPs mantenidas por la comunidad, como la \textit{Adversary Emulation Library} \cite{mitreAEL}, desarrollada por el \textit{MITRE Center for Threat-Informed Defense}. Este recurso proporciona perfiles de adversarios basados en comportamientos observados en grupos APT, aportando una base verificable para la creación de escenarios.

Esta información, estructurada mediante MITRE ATT\&CK y contextualizada a través de la \textit{Cyber Kill Chain}, permite diseñar simulaciones realistas, evaluar defensas, identificar brechas y entrenar al personal frente a amenazas alineadas con el panorama actual.

\section{Windows Internals}

El desarrollo y análisis de malware moderno está estrechamente vinculado al funcionamiento interno del sistema operativo. Comprender cómo Windows gestiona procesos, memoria, objetos del kernel y llamadas al sistema es fundamental, ya que constituye la base tanto de las técnicas ofensivas como defensivas. Los atacantes aprovechan estos mecanismos para ejecutar código de forma encubierta, mantener persistencia o evadir la detección, mientras que los analistas los utilizan para interpretar los rastros que dejan estas técnicas en el sistema.

Esta sección presenta los conceptos clave de la arquitectura interna de Windows relevantes para el análisis desarrollado en este trabajo y está pensada como material de referencia. Su objetivo es servir de apoyo a las explicaciones más técnicas que aparecen posteriormente, especialmente para lectores que puedan carecer de una base teórica sólida en sistemas operativos o internals de Windows.

\input{chapters/windows_internals/annex-virutual-memory.tex}
\newpage
\input{chapters/windows_internals/annex-process.tex}
\newpage
\input{chapters/windows_internals/annex-pe-files.tex}
\newpage
\subsection{DLL (Dynamic-Link Library)}

La mayoría de los ejecutables actuales dependen de librerías externas, algunas proporcionadas por el sistema operativo, como las funciones de entrada/salida de C (por ejemplo, \texttt{printf}, que se encuentra en \texttt{msvcrt.dll}), y otras que pueden ser librerías de terceros.  

Estas dependencias pueden manejarse de dos formas:

\begin{itemize}
    \item \textbf{Dependencias externas:} cargadas como \acrshort{dll} (\textit{Dynamic-Link Library}) en tiempo de ejecución.
    \item \textbf{Dependencias estáticas:} incluidas directamente en el binario final durante la compilación.
\end{itemize}

Las \acrshort{dll} son archivos con formato \acrshort{pe}, muy similares a los ejecutables. Cada módulo contiene sus propias cabeceras y secciones, tal como se explicó en el apartado de \acrshort{pe} files, y se mapean en memoria de manera similar a un ejecutable.

\subsubsection{Carga de DLLs}

Las \acrshort{dll} se cargan en el espacio de direcciones del proceso de dos formas:

\begin{itemize}
    \item \textbf{Carga estática:} Si el ejecutable indica durante la compilación que depende de una \acrshort{dll}, el \textit{loader} de Windows la cargará al iniciar el programa y resolverá la \textit{Import Address Table (\acrshort{iat})}.
    \item \textbf{Carga dinámica:} Si la \acrshort{dll} se carga en tiempo de ejecución mediante funciones como \texttt{LoadLibraryA}, el \textit{loader} mapea la librería en el espacio de direcciones y resuelve la \acrshort{iat} para el módulo que la cargó.
\end{itemize}

Una vez cargada, la \acrshort{dll} queda registrada en el \textit{Process Environment Block (\acrshort{peb})} del proceso. Si, en algún momento posterior, otro módulo intenta cargar la misma \acrshort{dll}, el sistema verifica el \acrshort{peb} y, en lugar de cargarla de nuevo, reutiliza la instancia ya existente en memoria.

\paragraph{Orden de búsqueda de DLLs en Windows}

El sistema operativo busca las \acrshort{dll} en el siguiente orden:

\begin{enumerate}
    \item Directorio desde el cual se cargó la aplicación
    \item \texttt{C:\textbackslash Windows\textbackslash System32}
    \item \texttt{C:\textbackslash Windows\textbackslash System}
    \item \texttt{C:\textbackslash Windows}
    \item Directorio actual de trabajo (Current Working Directory)
    \item Directorios listados en la variable de entorno \texttt{PATH} del sistema
    \item Directorios listados en la variable de entorno \texttt{PATH} del usuario
\end{enumerate}


Este mecanismo permite que un proceso utilice código externo sin incorporarlo de forma estática en su binario, favoreciendo la modularidad y la eficiencia en el uso de memoria.

Para el desarrollo del trabajo se realizó una implementación en C de una versión simplificada del loader de Windows. En la referencia \cite{walkingloader} se describen de forma detallada todos los pasos que ejecuta un programa en C para convertir el archivo \acrshort{pe} de una \acrshort{dll} en un módulo cargado en memoria y listo para ser utilizado.

\newpage
\input{chapters/windows_internals/annex-threads.tex}
\newpage
\input{chapters/windows_internals/annex-ring-levels.tex}
